\chapter{Starea actual\u a a domeniului}

\^In acest capitol este analizat contextul \^in care se \^incadreaz\u a aplica\c tia web pentru organizarea \c si gestionarea evenimentelor: cercet\u ari relevante pentru comer\c electronic, ticketing online \c si securitate web, platforme existente \c si o compara\c tie func\c tional\u a.

\section{Studii \c si articole \c stiin\c tifice}

\subsection{Arhitecturi web \c si API-uri REST}

Fielding \cite{fielding2000} define\c ste principiile REST (Representational State Transfer) ca stil arhitectural pentru sisteme distribuite. Separarea clientului de server, utilizarea resurselor adresabile prin URI \c si a metodelor HTTP standard (GET, POST, PUT, DELETE) faciliteaz\u a dezvoltarea aplica\c tiilor web scalabile. ManFast expune un API REST \^in Node.js/Express, consumat de frontend-ul React --- abordare aliniat\u a cu recomand\u arile din literatura de specialitate.

Fowler \cite{fowler2002} descrie tipare de proiectare enterprise, inclusiv separarea responsabilit\u a\c tilor \^intre straturi (prezentare, logic\u a de business, persisten\c t\u a). Aplica\c tia analizat\u a respect\u a acest principiu: React gestioneaz\u a UI, Express con\c tine regulile de business, PostgreSQL stocheaz\u a datele.

\subsection{Securitatea autentific\u arii}

OWASP \cite{owasp2024} enumer\u a bune practici pentru autentificare: hash-uirea parolelor (nu stocarea \^in clar), protec\c tia sesiunilor, limitarea ratei de cereri, validarea input-ului. ManFast implementeaz\u a bcrypt pentru parole, JWT cu expirare 24h, rate limiting pe codurile de verificare email \c si resetare parol\u a.

Standardul JWT (JSON Web Token) \cite{jwt2015} permite transmiterea securizat\u a a claims-urilor \^intre p\u ar\c ti. Token-ul semnat con\c tine id-ul utilizatorului \c si rolul (user/admin), verificat la fiecare cerere protejat\u a.

Provos \c si Mazi\`eres \cite{bcrypt1999} introduc bcrypt, algoritm de hash adaptat la cost computa\c tional configurabil, rezistent la atacuri brute-force pe parole. ManFast folose\c te cost factor 10.

\subsection{Pl\u a\c ti online \c si gateway-uri}

Documenta\c tia Stripe \cite{stripe2024} descrie fluxul Checkout: sesiune de plat\u a g\u azduit\u a, redirect c\u atre pagina securizat\u a Stripe, confirmare prin webhook sau verificare sesiune. ManFast creeaz\u a sesiuni Checkout cu metadata (event\_id, user\_id, nume participan\c ti) \c si confirm\u a plata prin endpoint dedicat, f\u ar\u a a stoca date de card pe serverul propriu --- reduc\^and responsabilitatea PCI-DSS.

\subsection{Persisten\c t\u a \c si baze de date rela\c tionale}

PostgreSQL \cite{postgresql2024} ofer\u a tranzac\c tii ACID, constr\u angeri de integritate referen\c tial\u a \c si tipuri avansate (ex.: array pentru imagini eveniment). Schema ManFast exploateaz\u a aceste capabilit\u a\c ti pentru consisten\c ta bilete--locuri disponibile.

\subsection{Interfe\c te utilizator \c si SPA}

React \cite{react2024} \c si ecosistemul s\u au permit construirea de Single Page Applications reactive. Material UI \cite{mui2024} ofer\u a componente accesibile \c si responsive. ManFast utilizeaz\u a React 19, Vite pentru build rapid \c si MUI 9 cu tem\u a dark personalizat\u a.

\section{Aplica\c tii dedicate domeniului}

\subsection{Eventbrite}

Eventbrite \cite{eventbrite2024} este o platform\u a interna\c tional\u a pentru crearea \c si promovarea evenimentelor, cu v\^anzare bilete, pagini de destina\c tie \c si instrumente de marketing. Ofer\u a check-in prin cod QR \c si analiz\u a audien\c t\u a. Platforma suport\u a multiple tipuri de evenimente: concerte, conferin\c te, workshop-uri. Fluxul tipic: organizatorul creeaz\u a pagina evenimentului, seteaz\u a pre\c turile \c si capacitatea, promoveaz\u a link-ul, iar participan\c tii achizi\c tioneaz\u a bilete online.

Limit\u ari pentru organizatori mici: comisioane per bilet, personalizare redus\u a a fluxului de check-in, dependen\c t\u a de ecosistemul proprietar. Datele participan\c tilor sunt stocate pe infrastructura Eventbrite, nu \^in controlul direct al organizatorului. Pentru conferin\c te cu sesiuni paralele sau concerte cu locuri numerotate, unele func\c tionalit\u a\c ti avansate pot lipsi sau necesita addon-uri pl\u a\c ite.

Din perspectiva tehnologic\u a, astfel de platforme folosesc arhitecturi cloud scalabile, procesare pl\u a\c ti prin PSP integrat \c si aplica\c tii mobile pentru scanare QR. ManFast adopt\u a un subset al acestor capabilit\u a\c ti, orientat spre simplitate \c si control local al datelor \^in PostgreSQL.

aici pun poza cu aplicatia 


\subsection{Platforme locale \c si generice}

Pe pia\c ta rom\^aneasc\u a exist\u a solu\c tii pentru bilete la concerte \c si festivaluri (iaBilet, Eventim), precum \c si platforme e-commerce generice (WooCommerce, Shopify) care pot vinde ``produse'' reprezent\^and locuri la eveniment. Acestea sunt optimizate pentru v\^anzare la scar\u a larg\u a, nu neap\u arat pentru meeting-uri corporate sau conferin\c te universitare.

Organizatorii locali folosesc frecvent formulare Google, pl\u a\c ti directe \c si liste Excel --- solu\c tii cu cost zero ini\c tial, dar cu risc ridicat de erori \c si lips\u a de audit trail pentru pl\u a\c ti. Reconcilierea manual\u a a listei de participan\c ti cu transferurile bancare consum\u a timp \c si poate genera dispute la intrarea \^in sal\u a.

Aplica\c tia propus\u a \^incerc\u a s\u a ofere un echilibru: cost de operare redus (self-hosted), flux clar de la descoperire eveniment la plat\u a \c si check-in, f\u ar\u a complexitatea unei platforme globale. Adaptarea la jude\c tele rom\^ane\c ti \c si pl\u a\c ti RON via Stripe reprezint\u a un avantaj competitiv \^in context local fa\c t\u a solu\c tiilor generice interna\c tionale.

\subsection{Tendințe și standarde în industria ticketing-ului}

Industria biletelor electronice a crescut accelerat după pandemia COVID-19, când contactless check-in și plățile online au devenit normă. Platformele mature investesc în trei direcții: (1) experiență mobilă — aplicații native sau PWA pentru scanare rapidă; (2) analitică — dashboard-uri de vânzări, heatmap-uri de prețuri dinamice; (3) integrare marketing — email automation, remarketing pe rețele sociale.

ManFast nu urmărește replicarea tuturor acestor capabilități enterprise, ci funcționalitatea minimă viabilă pentru organizatori locali: publicare eveniment, vânzare online, listă participanți, check-in la intrare. Această decizie de scope este aliniată cu principiul MVP (Minimum Viable Product) din literatura de product management.

\subsection{Conformitate PCI-DSS și delegarea plăților}

Standardul PCI-DSS (Payment Card Industry Data Security Standard) impune măsuri stricte pentru orice sistem care procesează, stochează sau transmite date de card. Soluțiile care folosesc Stripe Checkout hosted delegă colectarea datelor sensibile către Stripe — comerciantul primește doar confirmarea plății prin API. ManFast adoptă explicit acest model: frontend-ul redirecționează utilizatorul pe domeniul checkout.stripe.com, iar backend-ul verifică statusul sesiunii fără a vedea numărul de card.

Această abordare reduce suprafața de atac și simplifică auditul de securitate pentru o lucrare de licență, menținând totodată un flux profesional de plată recunoscut de utilizatori.

\subsection{Arhitectura SPA și API REST}

Aplicațiile Single Page Application (SPA) [8] încarcă o singură pagină HTML și navighează prin JavaScript, obținând date de la server prin API JSON. Avantaje: interfață reactivă, separare clară frontend/backend, posibilitatea de a extinde ulterior cu aplicație mobilă consumând același API.

Dezavantaje relevante pentru ManFast: SEO limitat pe pagini dinamice (acceptabil pentru platformă cu autentificare), necesitatea gestionării token-ului JWT pe client. Comparația cu arhitectura server-side rendering (Next.js) ar fi o extensie viitoare; pentru licență, Vite + React oferă timp de dezvoltare redus și ecosistem bogat (Material UI).

\subsection{Studiu de caz: fluxul utilizatorului în platforme consacrate}

Pe Eventbrite, utilizatorul parcurge: descoperire eveniment → selectare tip bilet → checkout → email confirmare cu PDF/QR. Organizatorul accesează dashboard cu statistici vânzări și export CSV participanți.

ManFast simplifică acest flux pentru organizatori locali (conferințe, concerte, meeting-uri):

-un singur tip de bilet pe eveniment (preț unic);

- nume nominal obligatoriu (identificare la check-in);

- fără marketplace global — evenimentele sunt listate doar pe instanța ManFast;

- check-in manual din panou admin, fără aplicație mobilă dedicată scanării QR (extensie viitoare).

Această simplificare reduce complexitatea implementării și timpul de învățare pentru administratori neexperimentați în IT.

\section{Criterii de evaluare a soluțiilor}

Pentru a fundamenta comparația din tabelul 2.2, au fost definite criterii ponderate:

Adecvare funcțională (40\%) — suport bilete nominale, check-in, filtrare locală.

Cost total de operare (20\%) — comisioane PSP, abonamente platformă, hosting.

Control date (20\%) — unde sunt stocate datele participanților, export, GDPR.

Personalizare (10\%) — posibilitatea modificării fluxului sau branding.

Securitate autentificare (10\%) — verificare email, reset parolă, roluri.

ManFast obține scor ridicat la control date și personalizare (cod deschis, PostgreSQL propriu), scor mediu la cost (Stripe per tranzacție, hosting self-managed), și scor bun la adecvare funcțională pentru nișa aleasă.

Aceste analize suplimentare contextualizează poziționarea ManFast față de soluțiile existente și justifică alegerile de scope ale proiectului.

\section{Analiz\u a comparativ\u a}

Tabelul~\ref{tab:comparatie} sintetizeaz\u a diferen\c tele \^intre platforme reprezentative \c si aplica\c tia propus\u a.

\begin{table}[htbp]
	\centering
	\begin{tabular}{|p{3.5cm}|p{2.2cm}|p{2.5cm}|p{3.5cm}|}
		\hline
		\textbf{Funcționalitate} & \textbf{Eventbrite} & \textbf{Generic e-comm} & \textbf{Aplicația propusă} \\ \hline
		Bilete nominale & Da & Parțial & Da \\ \hline
		Plăți Stripe RON & Nu (alte PSP) & Variabil & Da \\ \hline
		Check-in admin & Da (QR) & Nu & Da (listă) \\ \hline
		Verificare email register & Da & Variabil & Da (cod 6 cifre) \\ \hline
		Filtrare județ/dată & Da & Nu & Da (42 județe) \\ \hline
		Cod deschis / personalizabil & Nu & Nu & Da \\ \hline
		Reset parolă tel.+email & Variabil & Variabil & Da \\ \hline
	\end{tabular}
	\caption{Comparație funcțională}
	\label{tab:comparatie}
\end{table}

\section{Concluzii par\c tiale}

Analiza demonstreaz\u a c\u a exist\u a un spa\c tiu pentru o solu\c tie adaptat\u a: conferin\c te, meeting-uri, concerte \c si evenimente culturale, cu pl\u a\c ti RON via Stripe, bilete pe nume \c si administrare simpl\u a. Aplica\c tia propune o implementare open-source, extensibil\u a, care integreaz\u a func\c tionalit\u a\c tile esen\c tiale \c si le aliniaz\u a cu cerin\c tele locale (jude\c te, SMTP Brevo pentru emailuri).

\vspace{0.3cm}
\noindent Urm\u atorul capitol detaliaz\u a arhitectura, tehnologiile \c si implementarea solu\c tiei propuse.
